\begin{table}[!htbp]\centering\scriptsize\setlength{\tabcolsep}{4pt}
\caption{Ce qui fait varier le gain de la proposition et la neutralité du 1~\%}\label{tab:ircsg_sens}
\resizebox{\linewidth}{!}{%
\begin{tabular}{@{}lrl@{}}
\toprule
Hypothèse & Gain (Md€) & Autre effet \\
\midrule
Modèle, hypothèses centrales & 4,91 & participation 2,27~Md€ \\
Élasticités basses / hautes & 2,25 / 10,94 & \\
Élasticités de long terme (intensif 0,10 à 0,25 ; participation 0,10) & 7,30 & participation 0,91~Md€ \\
CSG perçue à moitié aujourd'hui & 1,29 & intensif $-$0,98~Md€ \\
CSG non perçue aujourd'hui & $-$2,09 & intensif $-$4,37~Md€ \\
Régime permanent : impôt futur sur les pensions vu à moitié par les actifs & 4,10 & \\
Régime permanent : vu en entier & 3,25 & coin $+$1,40~point \\
Répartition des revenus au PFU logés en haut (convention « haut ») & 4,5 & 1~\% $-$2,8~Md€ ; neutre si 50~\% à 110~k€ \\
Répartition des revenus au PFU au prorata & 5,3 & 1~\% +3,7~Md€ ; neutre si 50~\% à 858~k€ \\
\bottomrule
\end{tabular}}
\par\vspace{3pt}\parbox{\linewidth}{\raggedright\scriptsize\color{encre}
Gain : baisse de la perte sèche sur l'activité (jeu central sauf mention). CSG non perçue : le revenu d'activité ne réagit aujourd'hui qu'à l'IR (comme aux cotisations chez Lehmann, Marical et Rioux, 2013) et réagit après la fusion à tout l'impôt fusionné ; la marge de participation est gardée. Régime permanent : la hausse de l'impôt sur les pensions réduit la valeur des droits acquis par les cotisations contributives (27,4~\% du salaire brut) ; borne haute, le lien entre cotisation et pension étant faible au bas de l'échelle et nul au-delà de 8 plafonds. Convention centrale de répartition des revenus au PFU : « mixte » (controle\_hauts\_purs.py).}
\end{table}
